\section{Related Literature}
\label{sec:literature}

Cartels are a common practice and are quite costly to society. It is estimated that less than 15\% of US cartels are identified by the Department of Justice \citep{bryant1991price}. Considering that US anti-cartel policy is recognized as the most active, it is reasonable to assume that this proportion is even lower in countries with less antitrust enforcement tradition.

Cartels are associated with higher prices to final consumers \citep{connor2007price}, resulting in lower production than the socially optimum level and costs of maintaining the cartel and covering up its activities. For those reasons, the identification and subsequent punishment of cartels are a priority for virtually all competition authorities. Since it is a known illicit practice, it is not easy to detect cartels. There are two essential and complementary instruments to identify this anticompetitive behavior.

The first instrument is disclosure by one of the cartel's participants through leniency agreements or award-winning disclosure or of third-party claimants who had access to direct evidence of the agreement among the cartel's participants \citep{wils2007leniency}. Even if the competition authority develops a leniency program, its position is predominantly passive because third parties perform the identification of the cartel and the initial collection of evidence \citep{wils2016leniency}.

The second identification tool is predominantly active and consists of methods for analyzing companies' behavior in different markets with the aim of identifying suspicious behavior that is consistent with the existence of a cartel and inconsistent with the hypothesis of competition between companies. These methods are known as \textit{screening methods} \citep{harrington2008detecting}.

Observing suspicious behavior is not enough to prove a cartel's existence, but it is a crucial tool to guide competition authorities' investigations in their search for direct evidence. Evidence can be obtained through search and seizure operations or by intercepting communications, which are only authorized by the judiciary when there is robust evidence of suspicious conduct. Screening methods could help to identify this evidence.

Those two detection instruments---leniency agreements and screening methods---are complementary, not substitutes \citep{hüschelrath2014cartel,schinkel2013balancing}. Companies' incentives to opt for the leniency agreement by providing the competition authority with evidence of a cartel are associated with the probability of cartel detection. Proper screening methods can significantly increase the likelihood of identifying cartels.

In addition, the most stable cartels (i.e., those with a lower probability of detection) are possibly the most harmful to society and the least vulnerable to award-winning disclosure. For these reasons, authorities increasingly tend to use screening methods in their investigations \citep{schinkel2013balancing,imhof2018screening}. This trend is reinforced by increased data availability and processing capacity through machine learning and artificial intelligence, which give competition authorities a greater capacity to identify suspicious behaviors \citep{sanchezgraells2019screening}.

The literature on cartel detection models is reasonably prolific. In general, statistical models use observed price and quantity data to infer firms' patterns of behavior and reaction curves to identify whether such patterns are consistent with the cartel assumptions and inconsistent with the competition hypotheses.

Some of these models test whether the behavior of companies participating in a cartel differs from the action of nonparticipants, such as \citet{porter1993detection} and \citet{porter1999ohio}. In both cases, the model shows that the suspicion of cartel conduct, with defined participants and a defined period of operation, presents a correspondence with their market behavior, which has the value of proof if combined with a \textit{plus factor} \citep{hovenkamp2005federal}. However, it is not a full screening method; its role is not to identify suspicious conduct but to test the relevance of evidence that may have been brought by a complaint. Full screening methods use markers to signal which markets and firms are engaging in suspicious behavior.\footnote{A comprehensive review of screening models is presented by \citet{harrington2008detecting}.} These markers can be of three types.

The first deals with the relationship between firms' prices and demand movements, considering that a cartel structure reacts to demand fluctuations differently than firms that operate in competition. Some of the most notable papers on screening methods are from this first group, such as \citet{green1984noncooperative} and \citet{haltiwanger1991the}. The second group is about the relationship between market shares and price variance and is based on theoretical models and the empirical regularity of greater price stability in cartelized markets, which is broadly well documented in the literature \citep{abrantesmetz2006a,imhof2019detecting}. Finally, the third group focuses on the relationship between firms' prices, estimating reaction functions, and how they would be associated with competitive or collusive behavior \citep{bajari2003deciding}.

In the case of screening methods for public procurement, the variables used for cartel detection are specific to those observed in public tenders, such as the bidding pattern and the existence or absence of rotation among the winners, as is the standard of cartel operation; it has become conventional to call this \textit{bid-rigging} \citep{imhof2019detecting}. More recently, \citet{conley2016detecting} develop methods to detect bidders groups in collusive auctions, while \citet{chassang2022robust} propose robust screens that are less susceptible to strategic manipulation by cartels.

This pattern is essential in cases where the cartel participants do not have mechanisms for transferring the revenues derived from the cartel, so the rotation of winners becomes a necessary mechanism to avoid defections. However, this identification is not appropriate for cases where transfer mechanisms are feasible, either via corporate control or informally among cartel participants.

One of the difficulties that should be addressed by screening methods is their ability to rule out alternative hypotheses. For example, in the case of public bids, the bidding pattern may respond to cost differentials between companies, resulting in bids and ranking of winners that might appear to be the result of coordination between competitors.

An interesting proposal to circumvent this problem is offered by \citet{kawai2019using}, who use the margin of victory, analogous to a discontinuous regression method, to compare firms with supposedly similar cost structures to rule out this alternative hypothesis. They show that the inferences of the model or the robustness of the results is a recurring concern of the new screening methods.

There are two characteristics of the models mentioned that deserve to be highlighted because they emphasize the contributions of this paper. First, none of the cited models distinguishes tacit collusion from cases of explicit collusion, i.e., cartel \citep{harrington2008detecting}. However, this distinction is fundamental to the application of antitrust policy because, although the cartel is considered the most serious of antitrust offenses, tacit collusion is not even an offense and, therefore, not punishable.

The second characteristic is that the models mentioned above are a means of detecting a cartel after its occurrence. The variables used are prices, quantities, and the winners' rotation pattern observed in a given period. Therefore, these methods are here called \textit{ex-post screening models}. In conventional markets (i.e., those that are not public procurement), this qualification is unnecessary because the repression of cartels is typically carried out a posteriori, and preventive intervention is performed through the control of structures such as mergers and acquisitions \citep{hovenkamp2005federal}.

In the case of public procurement, however, the cartel is often held before the bidding process, i.e., before the participants disclose their bids. Identifying the cartel a posteriori fulfills the punitive nature by intending to inhibit future illicit conduct, but it does not prevent the loss of a fraudulent tender by the cartel's practice. These costs are exceptionally high in public bids, which require time for planning and high transaction costs to cancel contracts in execution and to conduct a new tender process.

Therefore, it would be desirable to develop \textit{ex-ante cartel screening methods} that could identify suspicious behavior before the execution of the bid. The use of frequent losers as a marker for collusion may consist of ex-ante screening, thereby distinguishing it from the other models in this regard. This is the main contribution of this research.

This marker can also be used in conjunction with other markers and variables in the ex-post analysis to strengthen the analysis, as recommended by \citet{harrington2008detecting}. The ``frequent losers'' marker might distinguish tacit collusion from explicit collusion, bringing more efficiency to the investigation techniques.\footnote{See Section~\ref{sec:results} for more details.}

Finally, it is relevant to note that cartelized companies respond strategically to the enforcement of antitrust authorities in what \citet{schinkel2013balancing} called a \textit{cat and mouse game}: cartelized companies will simultaneously seek to raise prices and to avoid detection by the competition authority. This situation may result in high prices coinciding with other market evidence that indicates increased competition, such as the number of participants and price dispersion.

In addition, cartel participants will seek to alter their behavior according to the detection model used to avoid being caught. For this reason, \citet{schinkel2013balancing} argues that the development of screening methods is a continuous process that is necessary every time cartelized companies change their behavior to avoid detection. Companies increase costs in this process, which may render the cartel unfeasible when avoiding detection is more costly than its benefits.
